%%
%%Implementatie van de oplossing
%%

\chapter{Implementatie van de oplossing}
\label{ch:ImplementatieVanDeOplossing}

\section{Verwerking van invoergegevens}

De automatiseringsflow start met het verzamelen van gegevens via een Microsoft Forms-formulier met de naam ``Onboarding''. Dit formulier wordt gebruikt door de personeelsdienst om de noodzakelijke gegevens van nieuwe medewerkers centraal aan te leveren aan de ICT-dienst.

Binnen het formulier worden verschillende gegevens opgevraagd die noodzakelijk zijn voor het uitvoeren van de onboarding. Hierbij gaat het onder andere om de voornaam, achternaam, startdatum, eventuele einddatum van het contract, functietitel, telefoonnummer, site en dienst van de medewerker. Daarnaast kunnen ook bijkomende configuraties meegegeven worden, zoals de gedeelde postvakken waarvoor toegang vereist is en de distributielijsten waaraan de gebruiker toegevoegd moet worden.

Om de kans op foutieve invoer te beperken, werd een extra controleveld toegevoegd waarbij bevestigd moet worden dat de ingevoerde gegevens nagekeken werden. Op deze manier wordt geprobeerd om de kwaliteit van de brongegevens zo hoog mogelijk te houden voordat de automatische verwerking start.

De toegang tot het formulier is beperkt via de beveiligingsinstellingen van Microsoft Forms. Enkel medewerkers van de personeelsdienst en de ICT-dienst kunnen het formulier invullen. In de praktijk wordt het formulier hoofdzakelijk gebruikt door de personeelsdienst, terwijl de ICT-dienst toegang behoudt voor testdoeleinden en voor het opnieuw uitvoeren van een onboardingproces wanneer een fout optreedt.

\section{Cloudflow binnen Power Automate}

Na het verzenden van het Microsoft Forms-formulier wordt automatisch een cloudflow binnen Microsoft Power Automate gestart. Deze flow vormt de centrale verwerkingslaag van de oplossing en staat in voor het ophalen, verwerken en voorbereiden van de ingevoerde gegevens.

De flow start met de trigger ``Wanneer er een nieuwe reactie wordt verzonden''. Deze trigger detecteert automatisch wanneer een nieuwe onboardingaanvraag werd ingediend via Microsoft Forms. Hierbij wordt het juiste formulier geselecteerd op basis van het unieke formulier-id.

Nadat de flow gestart werd, worden de details van de ingevulde aanvraag opgehaald via de stap ``Reactiedetails ophalen''. Deze stap maakt het mogelijk om alle ingevoerde waarden uit het formulier beschikbaar te maken binnen de flow. De verschillende antwoorden kunnen vervolgens gebruikt worden als dynamische variabelen binnen verdere stappen van het automatiseringsproces.

Omdat bepaalde velden binnen het formulier afhankelijk zijn van eerder gekozen antwoorden, moet de flow rekening houden met conditionele logica. Zo worden bepaalde diensten enkel weergegeven wanneer een specifieke site geselecteerd werd. Hierdoor verschillen de beschikbare invoervelden afhankelijk van de gekozen locatie.

Om deze gegevens uniform te kunnen verwerken binnen de verdere automatisering, worden ``Compose''-stappen gebruikt. Deze stappen zorgen ervoor dat de correcte dienst automatisch geselecteerd wordt op basis van de gekozen site binnen het formulier.

Codefragment ~\ref{lst:ComposeDepartement} vereenvoudigde expressie toont de gebruikte logica.

\begin{lstlisting}[
    caption={Samenstelling van de departementen},
    label={lst:ComposeDepartement}]
    if(site == "Gemeentehuis",
        dienst_gemeentehuis,
        if(site == "Onthaalpoort",
            dienst_onthaalpoort,
            if(site == "Cultuurhuis Stroming",
                dienst_cultuurhuis,
                andere_dienst)))
\end{lstlisting}

Door gebruik te maken van deze conditionele verwerking kan de flow steeds de correcte dienst ophalen, ongeacht welke site geselecteerd werd binnen het formulier. Deze waarden worden later gebruikt voor het bepalen van de correcte OU-structuur en groepsconfiguraties binnen Active Directory.

Naast het bepalen van de dienst wordt binnen de flow eveneens een rolcode samengesteld op basis van het departement van de gebruiker. Deze rolcodes worden later gebruikt binnen de PowerShell-scripts voor het toekennen van groepen, rechten en bijkomende configuraties.

Hiervoor wordt opnieuw gebruikgemaakt van een ``Compose''-stap waarin conditionele logica toegepast wordt. Afhankelijk van het gekozen departement wordt een specifieke rolcode toegekend.

Codefragment ~\ref{lst:ComposeRole} vereenvoudigde expressie toont een voorbeeld van deze verwerking.

\begin{lstlisting}[
    caption={Samenstelling van de rollen},
    label={lst:ComposeRole}]
    if(departement == "Informatica",
        "IT-user",
        if(departement == "Personeelsdienst",
            "HR-user",
            if(departement == "Financien",
                "FI-user",
                "user")))
\end{lstlisting}

\inputminted{powershell}{PowerAutomate.ps1}

Door gebruik te maken van gestandaardiseerde rolcodes kan de verdere verwerking binnen PowerShell en Active Directory eenvoudiger en consistenter verlopen. Daarnaast maakt deze aanpak het mogelijk om rechten en groepslidmaatschappen centraal te beheren op basis van functie of dienst.

Binnen de flow worden daarnaast verschillende variabelen samengesteld die later gebruikt worden binnen de PowerShell-scripts. Hierbij gaat het onder andere om de naam van de gebruiker, de dienst, de organisatorische eenheid binnen Active Directory en bijkomende configuraties zoals distributielijsten of gedeelde postvakken.

Tijdens de eerste ontwikkelingsfase werd gebruikgemaakt van een tijdelijke tussenstap waarbij de gegevens eerst naar een Excelbestand geschreven werden. Deze aanpak maakte het mogelijk om de verwerking van gegevens eenvoudig te controleren voordat de effectieve automatisering binnen Active Directory werd uitgevoerd.

Figuur ~\ref{fig:VoorbeeldOutput} toont voorbeeldoutput die gebruikt werd om te valideren dat de gegevens correct verwerkt werden.

\begin{figure}[htbp]
    \centering
    \includegraphics[width=1\textwidth]{Voorbeeldoutput.png}
    \caption[Voorbeeld output Power Automate]{Voorbeeld van de tijdelijke output van de Power Automate flow}
    \label{fig:VoorbeeldOutput}
\end{figure}

Figuur ~\ref{fig:Cloudflow_deel1} toont het eerste deel van de cloudflow binnen Microsoft Power Automate.

\begin{figure}[htbp]
    \centering
    \includegraphics[width=0.9\textwidth]{Cloudflow_deel1.png}
    \caption[Cloudflow binnen Power Automate]{Overzicht van het eerste deel van de cloudflow binnen Microsoft Power Automate.}
    \label{fig:Cloudflow_deel1}
\end{figure}

\section{Desktopflow en PowerShell-integratie}

Omdat Active Directory lokaal beheerd wordt binnen de infrastructuur van het gemeentebestuur, volstaat een gewone cloudflow binnen Microsoft Power Automate niet om gebruikersaccounts rechtstreeks aan te maken of aan te passen. De gebruikers dienen lokaal aangemaakt te worden aangezien er sprake is van een hybride omgeving. Hiervoor is communicatie nodig met de lokale serveromgeving waarop Active Directory beschikbaar is. 

Om deze communicatie mogelijk te maken, wordt gebruikgemaakt van Power Automate Desktop in combinatie met Power Automate Machine Runtime. Deze componenten vormen de verbinding tussen de cloudomgeving van Microsoft Power Automate en de lokale infrastructuur van de gemeente. Power Automate Desktop zal de scripts uitvoeren in een lokale flow en Power Automate Runtime zorgt voor de connectie tussen Power Automate Cloud en Power Automate Desktop. 

Wanneer de cloudflow alle gegevens verwerkt heeft, wordt automatisch een desktopflow gestart op de lokale server. Hierbij worden verschillende parameters doorgestuurd, waaronder de naam van de gebruiker, de dienst, de gekozen site, de rolcode en eventuele bijkomende configuraties zoals distributielijsten of gedeelde postvakken.

\subsection{Automatisch aanmaken en instellen van Active Directory-ge\-brui\-kers}

Binnen de desktopflow wordt een PowerShell-script uitgevoerd dat verantwoordelijk is voor het automatisch aanmaken en configureren van nieuwe gebruikersaccounts binnen Active Directory. Het script ontvangt verschillende parameters vanuit Power Automate, waaronder de voornaam, achternaam, dienst, site, functietitel en rolcode van de gebruiker.

Op basis van deze gegevens worden automatisch de correcte accountgegevens samengesteld, de juiste Organizational Unit bepaald en de noodzakelijke groepslidmaatschappen toegekend.

Codefragment ~\ref{lst:Parameters} toont de parameters die vanuit de desktopflow doorgestuurd worden naar het PowerShell-script.

\begin{lstlisting}[language=PowerShell,
    caption={Parameters},
    label={lst:Parameters}]
    param(
    [string]$FirstName,
    [string]$LastName,
    [string]$Department,
    [string]$Site,
    [string]$Role,
    [string]$Telephone,
    [string]$Title
    )
\end{lstlisting}

Door gebruik te maken van parameters kan hetzelfde script hergebruikt worden voor verschillende gebruikers en diensten zonder manuele aanpassingen aan de code.

\subsubsection{Verwerking van invoergegevens}

Voor het aanmaken van een account worden de ontvangen gegevens eerst gevalideerd en genormaliseerd. Speciale tekens en niet-alfabetische karakters worden verwijderd om problemen met accountnamen en e-mailadressen te vermijden. Door de tekst eerst te normaliseren, worden speciale tekens opgesplitst zodat deze eenvoudiger verwijderd kunnen worden.

Dezelfde verwerking wordt toegepast op zowel de voornaam als de achternaam van de gebruiker.

\begin{lstlisting}[language=PowerShell,
    caption={Aanpassen naam},
    label={lst:AanpassenNaam}]
    $FirstClean = $First.Normalize(
    [Text.NormalizationForm]::FormD
    ) -replace '\p{Mn}', ''
    
    $FirstClean = $FirstClean -replace '[^a-zA-Z]', ''
\end{lstlisting}

Deze verwerking zorgt ervoor dat gebruikersnamen consistent opgebouwd worden en compatibel blijven met de naamgevingsconventies binnen Active Directory.

Na het opschonen van de invoergegevens worden automatisch verschillende accountgegevens samengesteld. Hierbij gaat het onder andere om de \textit{SamAccountName}, het e-mailadres, de weergavenaam en bijkomende beschrijvende attributen.

\begin{lstlisting}[language=PowerShell,
    caption={Variabelen toekennen},
    label={lst:VariabelenToekennen}]
    $Sam = ($LastClean + $FirstClean.Substring(0,1)).ToLower()
    $Email = ($FirstClean + "." + $LastClean + "@berlare.be").ToLower()
    
    $DisplayName = "$First $Last"
    $Description = $Title
    $Company     = "Berlare"
\end{lstlisting}

Door deze waarden automatisch op te bouwen, wordt een uniforme naamgeving binnen de Active Directory-omgeving gegarandeerd.

Voordat een nieuwe gebruiker wordt aangemaakt, controleert het script eerst of reeds een account bestaat met dezelfde gebruikersnaam.

\begin{lstlisting}[
    caption={Dubbele controle},
    label={lst:DuplicateCheck}]{powershell}
    $Existing = Get-ADUser -Filter "SamAccountName -eq '$Sam'"
    
    if ($Existing) {
        Write-Warning "User already exists"
        return
    }
\end{lstlisting}

Hierdoor wordt vermeden dat dubbele accounts aangemaakt worden binnen Active Directory.

\subsubsection{Bepalen van de Organizational Unit}

Om gebruikers automatisch op de correcte locatie binnen Active Directory te plaatsen, worden eerst mogelijke OU-paden opgebouwd op basis van de gekozen site en dienst.

\begin{lstlisting}[
    language=PowerShell,
    caption={Samenstelling van de OU},
    label={lst:SamenstellingOU}]
    $BaseOU = "OU=Berlare Gemeente,DC=intranet,DC=berlare,DC=be"
    $UsersOUName = "OU=Users"
    
    $DeptOU = "$UsersOUName,OU=$Department,OU=Site $Site,$BaseOU"
    $SiteOU = "$UsersOUName,OU=Site $Site,$BaseOU"
\end{lstlisting}

Vervolgens controleert het script automatisch of een specifieke dienst-OU beschikbaar is. Indien deze niet gevonden wordt, valt het script terug op de algemene site-OU.

\begin{lstlisting}[language=PowerShell,
    caption={Controle van de OU},
    label={lst:ControleOU}]
    if ($DeptCheck) {
        $TargetOU = $DeptOU
    }
    else {
        $TargetOU = $SiteOU
    }
\end{lstlisting}

Deze aanpak zorgt ervoor dat het script flexibel kan omgaan met verschillen binnen de directorystructuur van de gemeente.

\subsubsection{Aanmaken van het gebruikersaccount}

Na het bepalen van de correcte OU wordt het gebruikersaccount effectief aangemaakt via het cmdlet \texttt{New-ADUser}. Hierbij worden verschillende attributen automatisch ingevuld, waaronder de naam, gebruikersnaam, functietitel, afdeling en het e-mailadres.

\begin{lstlisting}[language=PowerShell,
    caption={Creatie van de gebruiker},
    label={lst:UserCreation}]
    try {
        New-ADUser `
        -Name $DisplayName `
        -GivenName $First `
        -Surname $Last `
        -SamAccountName $Sam `
        -UserPrincipalName $Email `
        -Department $Department `
        -Title $Title `
        -Enabled $true `
        -Path $TargetOU
    }
    catch {
        Write-Error "Failed to create user $Sam"
        return
    }
\end{lstlisting}

Tijdens het aanmaken van de gebruiker wordt eveneens een tijdelijk wachtwoord ingesteld en wordt afgedwongen dat de gebruiker dit wachtwoord moet wijzigen bij de eerste aanmelding.

Door gebruik te maken van \texttt{try-catch}-constructies kunnen fouten gecontroleerd opgevangen en gelogd worden zonder dat het volledige script onverwacht stopt.

\subsubsection{Toekennen van groepslidmaatschappen}

Naast het aanmaken van het account worden automatisch verschillende groepslidmaatschappen toegekend. Hiervoor wordt gebruikgemaakt van vooraf gedefinieerde mappings tussen rolcodes en Active Directory-groepen.

\begin{lstlisting}[language=PowerShell,
    caption={Samenstelling rollen},
    label={lst:Rolemappings}]
    $RoleGroups = @{
        "IT-USER" = @("GG_GE_Informatica")
        "HR-USER" = @("GG_GE_Per")
        "FI-USER" = @("GG_GE_FINANCIEN")
    }
\end{lstlisting}

Daarnaast worden ook standaardgroepen toegevoegd die iedere gebruiker binnen de organisatie moet ontvangen.

\begin{lstlisting}[language=PowerShell,
    caption={Samenstelling groepen},
    label={lst:SamenstellingGroepen}]
    $DefaultGroups = @(
    "SEC_DR_Data2",
    "GG_GE_UNIFLOW",
    "SEC_DR_Fotoarchief"
    )
    
    $AllGroups = @()
    $AllGroups += $DefaultGroups
    
    foreach ($Group in $AllGroups) {
        Add-ADGroupMember -Identity $Group -Members $Sam
    }
\end{lstlisting}

Door groepslidmaatschappen automatisch toe te kennen, kunnen rechten en configuraties centraal beheerd worden via Active Directory-groepen.

\subsubsection{Terugkoppeling naar Power Automate}

Tot slot schrijft het script een beperkte uitvoer weg die opnieuw opgehaald kan worden binnen Power Automate. Hierbij wordt onder andere de User Principal Name teruggekoppeld naar de cloudflow.

\begin{lstlisting}[language=PowerShell,
    caption={Terugkoppeling UPN},
    label={lst:TerugkoppelingUPN}]
    $Result = [PSCustomObject]@{
        UPN = $Email
    }
    
    $Result | ConvertTo-Json -Compress
\end{lstlisting}

Deze terugkoppeling wordt later gebruikt binnen verdere automatiseringsstappen, zoals de synchronisatiecontrole met Entra ID en het toekennen van cloudlicenties.

\subsubsection{Instellen van accountvervaldatums}

Voor medewerkers met een tijdelijk contract moet het gebruikersaccount automatisch gedeactiveerd worden op het einde van de contractperiode. Om te vermijden dat deze vervaldatums manueel beheerd moeten worden, werd binnen de desktopflow een bijkomend PowerShell-script voorzien dat automatisch een accountvervaldatum instelt binnen Active Directory. Wanneer deze code binnen het algemene onboardingscript werd geplaatst, reageerde PowerShell hier niet op en stelde het de datum niet in. Een reden hiervoor werd niet gevonden. Wanneer deze code in een apart Power Automate blok werd gezet, werd de einddatum wel correct ingesteld. 

De einddatum wordt vanuit het Microsoft Forms-formulier doorgestuurd naar de desktopflow als parameter. Indien een einddatum ingevuld werd, verwerkt het script deze datum en koppelt deze aan het Active Directory-account van de gebruiker.

Codefragment ~\ref{lst:Einddatum} toont de verwerking van de meegegeven einddatum.

\begin{lstlisting}[language=PowerShell,
    caption={Verwerking van de mogelijke einddatum},
    label={lst:Einddatum}]
    $Sam = "%Username%"
    $DateString = "%ExpirationDate%"
    
    if (![string]::IsNullOrWhiteSpace($DateString)) {
        
        $DateObj = [datetime]::ParseExact(
        $DateString,
        "yyyy-MM-dd",
        $null
        )
        
        Set-ADAccountExpiration `
        -Identity $Sam `
        -DateTime $DateObj
    }
\end{lstlisting}

Het script controleert eerst of effectief een einddatum werd meegegeven. Indien dit niet het geval is, wordt geen vervaldatum ingesteld en blijft het account actief.\newline

Wanneer wel een einddatum aanwezig is, wordt deze geconverteerd naar een geldig \texttt{DateTime}-object zodat Active Directory deze correct kan verwerken. Vervolgens wordt via het cmdlet \texttt{Set-ADAccountExpiration} automatisch een accountvervaldatum ingesteld voor de gebruiker.

Binnen het uiteindelijke script werd bijkomende logica toegevoegd waarbij twee extra dagen toegevoegd worden aan de opgegeven einddatum. Daarnaast wordt het tijdstip ingesteld op 23:59:59 zodat het account actief blijft gedurende de volledige laatste werkdag van de medewerker.

\begin{lstlisting}[language=PowerShell,
    caption={Extra toevoegen 2 dagen bij einddatum},
    label={lst:Extra dagen}]
    $DateObj = $DateObj.AddDays(2)
    $DateObj = $DateObj.Date
    .AddHours(23)
    .AddMinutes(59)
    .AddSeconds(59)
\end{lstlisting}

Deze aanpak voorkomt dat accounts te vroeg gedeactiveerd worden door verschillen in synchronisatietijdstippen of tijdsinstellingen tussen systemen.

Ook binnen dit script wordt gebruikgemaakt van \texttt{try-catch}-constructies zodat fouten tijdens het instellen van de vervaldatum gecontroleerd opgevangen kunnen worden.

\begin{lstlisting}[language=PowerShell,
    caption={Try catch bij de einddatum},
    label={lst:TryCatchEinddatum}]
    try {
        Set-ADAccountExpiration `
        -Identity $Sam `
        -DateTime $DateObj
    }
    catch {
        Write-Output "ERROR:"
        Write-Output $_
    }
\end{lstlisting}

Door de vervaldatum automatisch te configureren, wordt de nood aan manuele opvolging verminderd en blijft het gebruikersbeheer consistenter en veiliger verlopen.

Figuur ~\ref{fig:DesktopFlow_deel1} toont het eerste deel van de desktopflow binnen Microsoft Power Automate.

\begin{figure}[htbp]
    \centering
    \includegraphics[width=1\textwidth]{DesktopFlow_deel1.png}
    \caption[Desktop flow deel 1]{Eerste flow op de desktop. In deze flow worden de gebruikers aangemaakt en variabelen terug gestuurd naar de cloud flow.}
    \label{fig:DesktopFlow_deel1}
\end{figure}

\section{Cloudverwerking binnen Microsoft 365 en Entra ID}

Deze sectie overloopt de uitwerking van de Entra ID en Microsoft 365 omgeving. 

\subsection{Synchronisatie met Entra ID}

Na het aanmaken van de gebruiker binnen de lokale Active Directory-omgeving moet het account eerst gesynchroniseerd worden naar Microsoft Entra ID voordat bijkomende cloudconfiguraties uitgevoerd kunnen worden. Binnen het gemeentebestuur wordt gebruikgemaakt van een hybride omgeving waarbij lokale Active Directory gekoppeld is aan Microsoft 365 via synchronisatieservices. Er wordt gebruik gemaakt van Entra ID Connect, vroeger ook wel bekend als Azure AD Connect. 

Tijdens deze synchronisatie worden gebruikersaccounts, attributen en groepsgegevens vanuit de lokale Active Directory doorgestuurd naar Entra ID. Hierdoor wordt het lokaal aangemaakte account ook beschikbaar binnen cloudtoepassingen zoals Microsoft 365, Exchange Online en Teams.

Omdat deze synchronisatie niet onmiddellijk gebeurt, kunnen cloudacties zoals het toekennen van licenties of Microsoft 365-groepen niet rechtstreeks uitgevoerd worden na het aanmaken van de gebruiker. De automatiseringsflow moet daarom eerst controleren of het account succesvol beschikbaar geworden is binnen Entra ID.

\subsection{Wachten op synchronisatie}

Om te controleren of de gebruiker succesvol gesynchroniseerd werd naar Entra ID, maakt de cloudflow gebruik van een \texttt{Do Until}-lus binnen Microsoft Power Automate. Deze lus blijft actief totdat het gebruikersaccount beschikbaar wordt binnen de cloudomgeving.

Voordat dit kan gebeuren, moet de aangemaakte UPN van de gebruiker doorgestuurd worden naar de cloudflow. Deze wordt automatisch doorgestuurd door de desktopflow, maar moet eerst nog verwerkt worden in de cloud. Dit gebeurd via de actie ``Parse JSON''. Deze zet de JSON waarde die doorgegeven wordt, via de connectie tussen cloud en desktop, om naar een variabele die normale tekst is. 

Daarna kan de controle starten. Bij het starten van deze controle wordt eerst een booleaanse variabele aangemaakt met de waarde \texttt{false}. Vervolgens probeert de flow via de actie ``Get user profile (V2)'' de gebruiker op te zoeken binnen Entra ID aan de hand van de eerder teruggekoppelde User Principal Name.

Wanneer het gebruikersprofiel succesvol opgehaald kan worden, wordt de variabele aangepast naar \texttt{true}, waardoor de \texttt{Do Until}-lus beëindigd wordt en de volgende cloudconfiguraties uitgevoerd kunnen worden.

Omdat de synchronisatie tussen Active Directory en Entra ID slechts één keer per half uur wordt gestart, werd binnen de lus eveneens een \texttt{Delay}-stap toegevoegd. Hierdoor wordt vermeden dat continu aanvragen naar Entra ID gestuurd worden. Binnen de implementatie wordt telkens 30 seconden gewacht voordat een nieuwe controle uitgevoerd wordt.

\subsection{Toekennen van Microsoft 365-groepen en licenties}

Nadat de gebruiker succesvol gesynchroniseerd werd naar Entra ID, kunnen bijkomende cloudconfiguraties uitgevoerd worden. Binnen de automatiseringsflow worden gebruikers automatisch toegevoegd aan vooraf gedefinieerde Microsoft 365-groepen.

Deze groepen worden gebruikt voor het centraal beheren van toegangsrechten en licenties binnen Microsoft 365. In plaats van licenties rechtstreeks toe te kennen aan individuele gebruikers, zijn de noodzakelijke licenties gekoppeld aan specifieke groepen binnen Entra ID. Wanneer een gebruiker toegevoegd wordt aan een dergelijke groep, worden de gekoppelde licenties automatisch toegewezen.

Binnen de implementatie worden gebruikers standaard toegevoegd aan een algemene Microsoft 365-groep voor E3-licenties die verantwoordelijk is voor de toekenning van basislicenties zoals Exchange Online en Microsoft Teams. Daarnaast wordt eveneens een bijkomende groep toegevoegd voor aanvullende functionaliteiten binnen de organisatie.

De gebruiker wordt naast de groep voor licenties van Microsoft 365 ook nog toegevoegd aan groepen voor een Power Automate licentie en een groep voor toegang tot de VPN. 

Door gebruik te maken van groepsgebaseerde licentietoekenning blijft het beheer eenvoudiger, consistenter en schaalbaarder. Nieuwe gebruikers ontvangen automatisch de correcte cloudconfiguraties zonder dat manuele tussenkomst vereist is.

Figuur ~\ref{fig:Do_until} toont de do until statement die in Power Automate gebruikt wordt.

\begin{figure}[htbp]
    \centering
    \includegraphics[width=1\textwidth]{Do_Until.png}
    \caption[Controle cloudgebruiker en toevoeging groepen]{Controle of de gebruiker aangemaakt is in Entra ID met de stappen om groepen toe te voegen.}
    \label{fig:Do_until}
\end{figure}

\section{Configuratie van e-mailtoegang}

Na het aanmaken van de gebruiker, de synchronisatie met Entra ID en het toekennen van licenties, worden bijkomende e-mailconfiguraties uitgevoerd. Hierbij gaat het voornamelijk om het toekennen van toegang tot gedeelde postvakken en het toevoegen van gebruikers aan distributiegroepen binnen Microsoft Exchange Online.

De benodigde configuraties worden geselecteerd binnen het Microsoft Forms-formulier tijdens de onboardingaanvraag. Deze gegevens worden vervolgens verwerkt binnen de cloudflow van Power Automate en doorgestuurd naar een bijkomende desktopflow waarin een PowerShell-script uitgevoerd wordt.

Binnen deze automatiseringsstap wordt gebruikgemaakt van de Exchange Online PowerShell-module om rechten toe te kennen op gedeelde postvakken en gebruikers automatisch toe te voegen aan de correcte distributiegroepen.

\subsection{Samenstellen van gedeelde postvakken en distributiegroepen}

Binnen de cloudflow worden de geselecteerde gedeelde postvakken en distributiegroepen opgehaald vanuit het Microsoft Forms-formulier. Omdat gebruikers meerdere opties kunnen selecteren, worden deze waarden verwerkt als lijsten die later doorgestuurd worden naar de desktopflow.

Bij het doorsturen van de informatie zonder ``Compose'' stappen, gaven ze niet altijd de correcte informatie door. Af en toe werden de arrays niet goed opgesteld. Hierdoor werden er toch ``Compose'' stappen toegevoegd zodat de gegevens telkens netjes in correcte arrays terecht kwamen. 

\subsection{Toekennen van mailboxrechten en distributiegroepen}

Omdat Power Automate Cloud geen stappen heeft voor het toekennen van rechten in Exchange Online, wordt er opnieuw gebruik gemaakt van Power Automate Desktop. Er wordt dus een nieuwe desktopflow opgesteld. 

In deze flow, wordt er gebruik gemaakt van de stap ``Toepassing uitvoeren''. De Exchange Online PowerShell module heeft inloggegevens nodig, of een certificaatgebaseerde authenticatie. Wegens tijdsgebrek is de certificaatgebaseerde authenticatie niet geconfigureerd. Dit deel van de oplossing is dus niet volledig geautomatiseerd. Er zullen dus inloggegevens ingevuld worden bij dit deel van de flow.

Het script ontvangt hierbij verschillende parameters vanuit Power Automate, waaronder de User Principal Name van de gebruiker, de geselecteerde gedeelde postvakken en de gekozen distributiegroepen.

Codefragment ~\ref{lst:ExchangeParameters} toont de parameters die vanuit de desktopflow doorgestuurd worden naar het script.

\begin{lstlisting}[language=PowerShell,
    caption={Parameters voor Exchange Online-configuratie},
    label={lst:ExchangeParameters}]
    param(
    [string]$UPN,
    [string]$Mailboxes,
    [string]$DistributionLists
    )
\end{lstlisting}

Voor de communicatie met Exchange Online wordt gebruikgemaakt van de module \texttt{ExchangeOnlineManagement}. Nadat verbinding gemaakt werd met Exchange Online, kan het script rechten en groepslidmaatschappen automatisch configureren.

\begin{lstlisting}[language=PowerShell,
    caption={Verbinding met Exchange Online},
    label={lst:ConnectExchange}]
    Import-Module ExchangeOnlineManagement
    
    Connect-ExchangeOnline -ShowBanner:$false
\end{lstlisting}

De geselecteerde gedeelde postvakken worden eerst verwerkt tot een lijst zodat meerdere mailboxen automatisch verwerkt kunnen worden binnen een lus.

\begin{lstlisting}[language=PowerShell,
    caption={Verwerking van mailboxlijsten},
    label={lst:MailboxList}]
    $MailboxList = $Mailboxes.Trim("[]") -split ","
    
    foreach ($sharedMailbox in $MailboxList) {
        
        $sharedMailbox = $sharedMailbox.Trim(
        " ","'","[","]"
        )
    }
\end{lstlisting}

Voor ieder gedeeld postvak wordt vervolgens automatisch \textit{Full Access}-toegang toegekend aan de gebruiker via het cmdlet \texttt{Add-MailboxPermission}.

\begin{lstlisting}[language=PowerShell,
    caption={Toekennen van mailboxrechten},
    label={lst:MailboxPermissions}]
    Add-MailboxPermission `
    -Identity $sharedMailbox `
    -User $UPN `
    -AccessRights FullAccess `
    -InheritanceType All `
    -Confirm:$false
\end{lstlisting}

Naast volledige toegang wordt eveneens automatisch \textit{Send on Behalf}-toegang ingesteld. Omdat een mailbox hiervoor verplicht is, moet de configuratie van de E3 licentie afgerond zijn. Als deze afgerond is, kunnen de correcte rechten toegewezen worden. Er wordt hierbij gebruikgemaakt van een retry-mechanisme waarbij meerdere pogingen uitgevoerd worden met een tussentijdse wachttijd. Hierdoor worden, op het moment dat de mailbox aangemaakt is, de correcte mailboxen toegewezen aan de gebruiker.

\begin{lstlisting}[language=PowerShell,
    caption={Retry-mechanisme voor Send on Behalf},
    label={lst:RetrySendOnBehalf}]
    while (-not $Success -and $Attempt -lt $MaxAttempts) {
        
        try {
            
            Set-Mailbox `
            -Identity $sharedMailbox `
            -GrantSendOnBehalfTo @{Add=$UPN}
            
            $Success = $true
        }
        catch {
            
            $Attempt++
            
            Start-Sleep -Seconds 60
        }
    }
\end{lstlisting}

\subsection{Toevoegen aan distributiegroepen}

Naast gedeelde postvakken worden gebruikers eveneens automatisch toegevoegd aan distributiegroepen binnen Exchange Online. Hierbij worden zowel standaard distributiegroepen als bijkomende groepen uit het formulier verwerkt.

\begin{lstlisting}[language=PowerShell,
    caption={Standaard distributiegroepen},
    label={lst:DefaultDLs}]
    $DefaultDLs = @(
    "iedereen@berlare.be",
    "viruswaarschuwing@berlare.be"
    )
\end{lstlisting}

Tot slot worden alle distributiegroepen doorlopen en wordt de gebruiker automatisch toegevoegd via het cmdlet \texttt{Add-DistributionGroupMember}.

\begin{lstlisting}[language=PowerShell,
    caption={Toevoegen aan distributiegroepen},
    label={lst:AddDistributionGroup}]
    Add-DistributionGroupMember `
    -Identity $dl `
    -Member $UPN `
    -BypassSecurityGroupManagerCheck
\end{lstlisting}

Door deze configuraties automatisch uit te voeren, wordt het manuele beheer van mailboxrechten en distributiegroepen sterk verminderd en verloopt de onboarding consistenter en efficiënter.

\section{Foutafhandeling en betrouwbaarheid}

Binnen de verschillende PowerShell-scripts wordt gebruikgemaakt van foutafhandeling om onverwachte problemen gecontroleerd op te vangen zonder dat het volledige onboardingproces abrupt stopt. Hiervoor worden voornamelijk \texttt{try-catch}-constructies gebruikt.

Wanneer een fout optreedt tijdens het uitvoeren van een opdracht, wordt een waarschuwing of foutmelding weggeschreven zodat deze later geraadpleegd kan worden binnen Power Automate of tijdens troubleshooting. Dit zorgt er echter niet voor dat de flow automatisch stopt.

Codefragment ~\ref{lst:MailboxErrorHandling} toont een voorbeeld van foutafhandeling tijdens het toevoegen van een gebruiker aan een gedeeld postvak.

\begin{lstlisting}[language=PowerShell,
    caption={Foutafhandeling bij mailboxrechten},
    label={lst:MailboxErrorHandling}]
    try {
        
        Add-MailboxPermission `
        -Identity $sharedMailbox `
        -User $UPN `
        -AccessRights FullAccess `
        -ErrorAction Stop
    }
    catch {
        
        Write-Warning "Failed FullAccess"
        Write-Warning $_.Exception.Message
    }
\end{lstlisting}

Daarnaast wordt binnen bepaalde onderdelen gebruikgemaakt van retry-mechanismen. Dit is voornamelijk noodzakelijk bij synchronisatieprocessen binnen Exchange Online en Entra ID, waarbij objecten niet altijd onmiddellijk beschikbaar zijn na het aanmaken van een gebruiker.

Door meerdere pogingen uit te voeren met een tussentijdse wachttijd, kan het script tijdelijke synchronisatieproblemen automatisch opvangen zonder manuele tussenkomst.

\begin{lstlisting}[language=PowerShell,
    caption={Retry-mechanisme},
    label={lst:RetryMechanism}]
    while (-not $Success -and $Attempt -lt $MaxAttempts) {
        
        try {
            
            Set-Mailbox `
            -Identity $sharedMailbox `
            -GrantSendOnBehalfTo @{Add=$UPN}
            
            $Success = $true
        }
        catch {
            
            $Attempt++
            
            Start-Sleep -Seconds 60
        }
    }
\end{lstlisting}

Deze aanpak verhoogt de betrouwbaarheid van het onboardingproces en vermindert de nood aan manuele correcties wanneer tijdelijke fouten of synchronisatievertragingen optreden.

\section{Automatische terugkoppeling via e-mail}

Na het succesvol uitvoeren van de onboardingflow wordt automatisch een e-mail verstuurd vanuit een gedeeld postvak van de ICT-dienst. Deze e-mail bevat bijkomende informatie en handleidingen voor de nieuwe medewerker.

In deze handleidingen zijn stappenplannen te vinden met betrekking tot het gebruik van de VPN, het vrijgeven van e-mails die tegengehouden worden door de spamfilter en printen. In de e-mail zelf zijn ook een paar praktische afspraken in verband met het gebruik van wifi. Ook de uitleg over MFA is te vinden in deze e-mail. Ook is er in de e-mail een link naar de helpdesk waar ze al hun problemen moeten melden en uiteraard ook geholpen worden indien nodig. 

Binnen Power Automate moeten deze bijlagen eerst geladen worden. Dit gebeurd met de stap ``Get file content''. Voor de 4 bijlagen gebeurt dit dus op identieke, parallelle wijze. Eens deze opgeladen zijn in Power Automate, kunnen ze geselecteerd worden in de stap ``Send an email from a shared mailbox''. Hierin staat het e-mailadres dat de e-mail verstuurt, als ook de ontvanger. Dit is op dynamische wijze opnieuw de UPN variabele die initieel via de desktopflow naar de cloud wordt gestuurd. Net zoals bij een normale e-mail hoort er ook een onderwerp en de gehele e-mail bij. 

Figuur ~\ref{fig:Mails} toont hoe de bijlagen worden aangemaakt en hoe de e-mail verstuurd wordt.

\begin{figure}[htbp]
    \centering
    \includegraphics[width=1\textwidth]{Mails.png}
    \caption[Stappen om een e-mail te versturen]{Stappen die nodig zijn in een flow om een automatische e-mail met bijlagen te versturen}
    \label{fig:Mails}
\end{figure}

\section{Logging van uitgevoerde onboardings}

Om opvolging en controle mogelijk te maken, wordt na een succesvolle onboarding automatisch een logregistratie toegevoegd binnen een Excelbestand. Deze registratie bevat onder andere de naam van de gebruiker, de uitvoeringsdatum en de status van de onboarding.

Binnen de cloudflow wordt hiervoor gebruikgemaakt van de actie ``Update a row''. Hiermee wordt de bestaande aanvraag bijgewerkt zodat zichtbaar wordt dat de onboarding succesvol uitgevoerd werd.

Door gebruik te maken van deze logging kan de ICT-dienst eenvoudig controleren welke onboardings reeds uitgevoerd werden en kunnen eventuele problemen later sneller opgespoord worden. Ook voor de medewerkers van andere diensten is dit handig om te zien of alles van IT in orde is voor de nieuwe collega.

Figuur ~\ref{fig:Voorbeeld_output_logging} toont het loggen wanneer de gebruiker aangemaakt is.

\begin{figure}[htbp]
    \centering
    \includegraphics[width=1\textwidth]{VoorbeeldoutputLogging.png}
    \caption[Voorbeeld output van logging]{Aangepaste rij nadat de gebruiker is aangemaakt. Kolom ``Processed'' is veranderd van ``No'' naar ``Yes''}
    \label{fig:Voorbeeld_output_logging}
\end{figure}

\section{Geautomatiseerde offboarding van gebruikers}

Naast het onboardingproces werd eveneens een geautomatiseerde offboardingflow ontwikkeld binnen Microsoft Power Automate. Deze flow zorgt ervoor dat gebruikersaccounts en bijhorende configuraties automatisch aangepast worden wanneer een medewerker de organisatie verlaat.

De offboarding start vanuit een afzonderlijk Microsoft Forms-formulier waarin onder andere de User Principal Name van de gebruiker, de einddatum van het contract en de reden van uitdiensttreding ingevuld worden. Nadat het formulier verzonden werd, start automatisch een cloudflow binnen Power Automate die de verdere verwerking uitvoert.

Net zoals bij de onboarding wordt gebruikgemaakt van een hybride aanpak waarbij cloudflows gecombineerd worden met desktopflows en PowerShell-scripts. Hierdoor kunnen zowel cloudconfiguraties binnen Microsoft 365 als lokale configuraties binnen Active Directory automatisch aangepast worden.

Tijdens de offboarding worden verschillende acties uitgevoerd. Hierbij wordt de mailbox van de gebruiker eerst omgezet naar een gedeeld postvak waarna rechten op gedeelde postvakken en distributiegroepen verwijderd worden. Vervolgens wordt het Active Directory-account gedeactiveerd, verborgen uit de globale adreslijst en verwijderd uit beveiligingsgroepen en andere groepslidmaatschappen.

Daarnaast worden ook Microsoft 365-groepen aangepast zodat gekoppelde licenties automatisch vrijgegeven worden. Door deze stappen automatisch uit te voeren, verloopt de offboarding consistenter en wordt de nood aan manuele tussenkomst verminderd.

\subsection{Cloudflow voor offboarding}

De offboardingflow start binnen Microsoft Power Automate wanneer een nieuwe reactie verzonden wordt via het Microsoft Forms-formulier voor offboarding. Net zoals bij de onboarding worden de ingevulde gegevens automatisch opgehaald aan de hand van het unieke reactie-id van het formulier.

Na het ophalen van de gegevens wordt gecontroleerd welke einddatum voor de medewerker werd opgegeven. Deze datum bepaalt wanneer de effectieve offboarding uitgevoerd mag worden. Zolang de einddatum nog niet bereikt werd, blijft de flow in wachtstatus staan zodat de verdere automatiseringsstappen nog niet uitgevoerd worden.

Deze aanpak voorkomt dat een gebruikersaccount te vroeg gedeactiveerd wordt terwijl de medewerker nog actief is binnen de organisatie. Aangezien de personeelsdienst de einddatum reeds kent bij het registreren van de uitdiensttreding, kan deze informatie vooraf meegegeven worden binnen het formulier.

Binnen de cloudflow wordt hiervoor gebruikgemaakt van een wachttijd op basis van datum en tijd. Zodra de opgegeven einddatum bereikt wordt, hervat de flow automatisch en worden de verdere offboardingstappen uitgevoerd.

\subsection{Desktopflow en verwerking binnen Exchange Online}

Wanneer de opgegeven einddatum bereikt werd, start de cloudflow automatisch een desktopflow binnen Power Automate Desktop. Deze desktopflow is verantwoordelijk voor het uitvoeren van de PowerShell-scripts die de effectieve offboarding van de gebruiker uitvoeren binnen Exchange Online.

Net zoals bij de onboarding wordt gebruikgemaakt van een hybride aanpak waarbij Power Automate Cloud communiceert met de lokale infrastructuur via Power Automate Desktop en Power Automate Machine Runtime. Hierdoor kunnen lokale Active Directory-configuraties gecombineerd worden met cloudconfiguraties binnen Microsoft 365.

Vanuit de cloudflow worden verschillende parameters doorgestuurd naar de desktopflow, waaronder de User Principal Name van de gebruiker, de reden van uitdiensttreding en de einddatum van het contract.

Codefragment ~\ref{lst:OffboardingExchangeParameters} toont de parameters die gebruikt worden binnen het Exchange Online-script.

\begin{lstlisting}[language=PowerShell,
    caption={Parameters voor Exchange Online-offboarding},
    label={lst:OffboardingExchangeParameters}]
    param(
    [string]$UPN
    )
\end{lstlisting}

Binnen het script wordt eerst verbinding gemaakt met Exchange Online via de module \texttt{ExchangeOnlineManagement}. Hierdoor kunnen mailboxen, rechten en distributiegroepen automatisch aangepast worden tijdens de offboarding.

\begin{lstlisting}[language=PowerShell,
    caption={Verbinding met Exchange Online},
    label={lst:ConnectExchangeOffboarding}]
    Import-Module ExchangeOnlineManagement
    
    Connect-ExchangeOnline
\end{lstlisting}

\subsubsection{Omzetten naar gedeeld postvak}

De mailbox van de gebruiker wordt eerst automatisch omgezet naar een gedeeld postvak. Hierdoor blijven bestaande e-mails en agenda-items behouden nadat de gebruiker uit dienst gegaan is.

\begin{lstlisting}[language=PowerShell,
    caption={Omzetten naar gedeeld postvak},
    label={lst:ConvertSharedMailbox}]
    Set-Mailbox `
    -Identity $UserMailbox.PrimarySmtpAddress `
    -Type Shared
\end{lstlisting}

Door de mailbox om te zetten naar een gedeeld postvak kan de Microsoft 365-licentie later verwijderd worden zonder dat historische communicatie verloren gaat. Dit maakt het mogelijk om oude e-mails nog tijdelijk te raadplegen indien nodig.

\subsubsection{Verwijderen van mailboxrechten en distributiegroepen}

Na het omzetten van de mailbox worden bestaande rechten en groepslidmaatschappen verwijderd. Hierbij wordt eerst gecontroleerd op mailboxen waarvoor de gebruiker \textit{Full Access}-rechten bezit.

\begin{lstlisting}[language=PowerShell,
    caption={Verwijderen van Full Access-rechten},
    label={lst:RemoveFullAccess}]
    Remove-MailboxPermission `
    -Identity $Mailbox.PrimarySmtpAddress `
    -User $UPN `
    -AccessRights FullAccess `
    -Confirm:$false
\end{lstlisting}

Daarnaast wordt eveneens gecontroleerd of de gebruiker \textit{Send on Behalf}-rechten bezit op andere mailboxen. Indien deze rechten aanwezig zijn, worden ze automatisch verwijderd.

\begin{lstlisting}[language=PowerShell,
    caption={Verwijderen van Send on Behalf},
    label={lst:RemoveSendOnBehalf}]
    Set-Mailbox `
    -Identity $Mailbox.PrimarySmtpAddress `
    -GrantSendOnBehalfTo @{
        Remove = $UserMailbox.DistinguishedName
    }
\end{lstlisting}

Tot slot wordt de gebruiker automatisch verwijderd uit distributiegroepen binnen Exchange Online.

\begin{lstlisting}[language=PowerShell,
    caption={Verwijderen uit distributiegroepen},
    label={lst:RemoveDistributionGroups}]
    Remove-DistributionGroupMember `
    -Identity $_.PrimarySmtpAddress `
    -Member $UPN `
    -Confirm:$false
\end{lstlisting}

Door deze rechten en groepslidmaatschappen automatisch te verwijderen, wordt vermeden dat voormalige medewerkers nog toegang behouden tot gedeelde mailboxen of interne communicatiestromen.

\subsection{Desktopflow en verwerking binnen Active Directory}

Naast de configuraties binnen Exchange Online worden tijdens de offboarding eveneens verschillende aanpassingen uitgevoerd binnen de lokale Active Directory-omgeving. Hiervoor wordt een tweede PowerShell-script uitgevoerd vanuit Power Automate Desktop.

Dit script ontvangt verschillende parameters vanuit de cloudflow, waaronder de User Principal Name van de gebruiker, de reden van uitdiensttreding en de einddatum van het contract.

Codefragment ~\ref{lst:ADOffboardingParameters} toont de parameters die gebruikt worden binnen het Active Directory-script.

\begin{lstlisting}[language=PowerShell,
    caption={Parameters voor Active Directory-offboarding},
    label={lst:ADOffboardingParameters}]
    param(
    [string]$UPN,
    [string]$Reason,
    [string]$EndDate
    )
\end{lstlisting}

Binnen het script wordt eerst gecontroleerd of het gebruikersaccount effectief bestaat binnen Active Directory. Hiervoor wordt de gebruiker opgezocht aan de hand van de User Principal Name.

\begin{lstlisting}[language=PowerShell,
    caption={Opzoeken van de gebruiker},
    label={lst:GetADUserOffboarding}]
    $User = Get-ADUser `
    -Filter "UserPrincipalName -eq '$UPN'" `
    -Properties MemberOf,Description
\end{lstlisting}

Na het ophalen van het account kunnen de verdere offboardingstappen automatisch uitgevoerd worden.

\subsubsection{Deactiveren van gebruikersaccounts}

Tijdens de offboarding wordt het gebruikersaccount automatisch gedeactiveerd binnen Active Directory. Hierdoor kan de voormalige medewerker niet langer aanmelden binnen de omgeving van de organisatie.

\begin{lstlisting}[language=PowerShell,
    caption={Deactiveren van gebruikersaccounts},
    label={lst:DisableADAccount}]
    Disable-ADAccount -Identity $User
\end{lstlisting}

Daarnaast wordt het account verborgen uit de Global Address List van Exchange. Hierdoor verschijnt de gebruiker niet langer binnen adreslijsten van Microsoft Outlook of andere Microsoft 365-toepassingen.

\begin{lstlisting}[language=PowerShell,
    caption={Verbergen uit de Global Address List},
    label={lst:HideFromGAL}]
    Set-ADUser `
    -Identity $User `
    -Replace @{
        msExchHideFromAddressLists=$true
    }
\end{lstlisting}

Naast het uitschakelen van het account wordt eveneens de beschrijving van de gebruiker aangepast. Hierbij wordt automatisch geregistreerd vanaf wanneer de medewerker uit dienst is en wat de reden van uitdiensttreding was.

\begin{lstlisting}[language=PowerShell,
    caption={Aanpassen van de beschrijving},
    label={lst:DescriptionOffboarding}]
    $NewDescription = 
    "Uit dienst sinds $FormattedDate - $Reason"
    
    Set-ADUser `
    -Identity $User `
    -Description $NewDescription
\end{lstlisting}

Door deze informatie rechtstreeks binnen Active Directory op te slaan, blijft bijkomende administratieve informatie beschikbaar voor de ICT-dienst en systeembeheerders.

\subsubsection{Verwijderen van groepslidmaatschappen}

Na het deactiveren van het account worden bestaande groepslidmaatschappen automatisch verwijderd. Hierdoor verliest de gebruiker alle gekoppelde rechten en toegangen binnen de infrastructuur van de organisatie.

Binnen het script wordt eerst een beperkte lijst samengesteld van groepen die behouden moeten blijven. In de huidige implementatie blijft enkel de standaardgroep \texttt{Domain Users} behouden.

\begin{lstlisting}[language=PowerShell,
    caption={Groepen die behouden blijven},
    label={lst:GroupsToKeep}]
    $GroupsToKeep = @(
    "Domain Users"
    )
\end{lstlisting}

Vervolgens worden alle groepslidmaatschappen van de gebruiker opgehaald en gecontroleerd. Indien een groep niet voorkomt binnen de uitzonderingslijst, wordt de gebruiker automatisch uit deze groep verwijderd.

\begin{lstlisting}[language=PowerShell,
    caption={Verwijderen van groepslidmaatschappen},
    label={lst:RemoveADGroups}]
    Remove-ADGroupMember `
    -Identity $_ `
    -Members $User `
    -Confirm:$false
\end{lstlisting}

Door groepslidmaatschappen automatisch te verwijderen, wordt vermeden dat voormalige medewerkers nog toegang behouden tot netwerkshares, beveiligingsgroepen of andere interne systemen die gekoppeld zijn aan Active Directory-groepen.

Net zoals bij de onboarding wordt ook binnen deze scripts gebruikgemaakt van \texttt{try-catch}-constructies zodat fouten gecontroleerd opgevangen kunnen worden zonder dat het volledige proces abrupt stopt.

Figuur ~\ref{fig:DesktopOffboarding} toont de complete desktopflow voor het offboardingsproces.

\begin{figure}[htbp]
    \centering
    \includegraphics[width=1\textwidth]{DesktopOffboarding.png}
    \caption[Desktopflow offboarding]{Desktopflow voor het offboardingsproces}
    \label{fig:DesktopOffboarding}
\end{figure}

\subsection{Verwijderen van groepslidmaatschappen en licenties}

Naast de lokale Active Directory-groepen moeten tijdens de offboarding eveneens verschillende cloudgroepen binnen Microsoft 365 verwijderd worden. Deze groepen werden tijdens de onboarding automatisch toegekend via Power Automate en zijn verantwoordelijk voor het toewijzen van licenties en bijkomende toegangsrechten.

Binnen de implementatie gaat het onder andere om de Microsoft 365 E3-licentie, een Power Automate-licentie en een groep die toegang verleent tot de VPN-omgeving van de organisatie.

Het verwijderen van deze groepen gebeurt volledig binnen de cloudflow van Power Automate. Hiervoor wordt gebruikgemaakt van de actie ``Lid uit groep verwijderen''. Via deze stap wordt de gebruiker automatisch verwijderd uit de vooraf gedefinieerde Microsoft 365-groepen.

Voor het identificeren van de gebruiker wordt het e-mailadres gebruikt dat eerder via het formulier en de desktopflow verwerkt werd. Hierdoor kan de correcte gebruiker automatisch opgezocht en verwijderd worden uit de overeenkomstige groepen binnen Entra ID.

Door gebruikers automatisch uit deze groepen te verwijderen, worden gekoppelde licenties en toegangsrechten eveneens automatisch ingetrokken. Dit vermindert het risico op ongewenste toegang na uitdiensttreding en zorgt ervoor dat ongebruikte licenties opnieuw beschikbaar worden binnen de organisatie. Dit is dus zowel positief voor veiligheid van gegevens, als ook voor het niet verspillen van licenties. \newline

Figuur ~\ref{fig:CloudOffboarding} toont de complete cloudflow voor het offboardingsproces.

\begin{figure}[htbp]
    \centering
    \includegraphics[width=0.8\textwidth]{CloudOffboarding.png}
    \caption[Cloudflow offboarding]{Cloudflow voor het offboardingsproces}
    \label{fig:CloudOffboarding}
\end{figure}

%Om een automatische werking op te stellen, werd door een collega gekozen voor Power Automate. Power Automate is een programma van Microsoft zelf dat gebruikt kan worden om automatische flows te maken. Power automate is voornamelijk een tool dat in de cloud draait, maar het kan ook voor automatisatie zorgen on prem. \newline

%Voor cloud toepassingen van Microsoft heeft Power Automate verschillende stappen. Zo kan je echt op vlak van gebruikers mensen gaan aanmaken in Entra ID en deze aan groepen toevoegen etc. Aangezien er in de gemeente gebruik wordt gemaakt van een hybride omgeving, met dus zowel toepassing van Active Directory in de cloud en on prem, is het dus belangrijk dat de gebruikers on prem worden aangemaakt. \newline

%Voor active directory heeft Power Automate geen speciale stappen die je kan gebruiken. Hiervoor wordt het script, dat in hoofdstuk ~\ref{ch:analyseontwerp} uitgelegd wordt, gebruikt. Dit script staat dus lokaal op een server opgeslagen. Aangezien Power Automate een cloud tool is, is er de desktop versie van Power Automate nodig. Deze applicatie werkt als brug tussen de online versie en de lokale toepassing van de Active Directory. \newline

%In de cloud is er een flow gemaakt. Deze flow start met het ophalen van gegevens wanneer er een nieuwe inzending is van het formulier. Hiervoor is een stap die voorzien is, namelijk de stap ``Wanneer er een nieuwe reactie wordt verzonden''. In deze stap kan je eenvoudig in een keuze menu kiezen voor het formulier dat je in Microsoft Forms hebt gemaakt. Als deze er niet bijstaat, kan je ook op zoek gaan naar het id van het formulier. Deze kan je vinden in de URL van jouw form, achter het duidelijke ``FormID'' in de URL. \newline

%Hierna moet er nog een stap komen met de gemaakte form. Hiervoor heb je de stap ``Reactiedetails ophalen''. Hierin moet je opnieuw het id van de form ingeven, als ook het id van het antwoord kiezen. Bij het antwoord moet je kiezen voor een dynamische expressie te kiezen. Op deze manier kan je de laatst verzonden input selecteren. Je kan de dynamische expressie kiezen door aan de rechterkant van de kolom op de bliksemschicht te klikken, of door een ``/'' te typen.\newline

%De gemeente heeft, zoals uitgelegd in hoofdstuk ~\ref{ch:analyseontwerp}, op verschillende niveaus gebruikers. Ze zitten allemaal in een ``OU'' ``Users'', maar ze zitten wel nog steeds, afhankelijk van de sites, 1 niveau verschil. Hiervoor in binnen Power Automate, een extra stap gemaakt om de juiste sites en diensten toe te kennen. Hiervoor is de stap ``Compose'' aangeroepen. 

%Deze ``Compose'' gaat bij de vraag van welke site, kijken naar de reactie. Op vlak daarvan gaat hij het antwoord van de vraag over welke dienst de gebruiker zit. Zo gaat de form bijvoorbeeld bij het selecteren van site ``Gemeentehuis'', naar de vraag voor welke diensten van site gemeentehuis. Hierin staan dan bijvoorbeeld diensten ``vergunningen, personeelsdienst of informatica''. In de IF-statement kijkt hij bij het selecteren dan naar die vraag. Dit zorgt er voor dat hij dus altijd voor de dienst, de correcte dienst die in de form wordt gekozen, wordt ingevuld. Om het heel simpel uit te leggen, de IF-statements simuleren de werking van het formulier dat over bepaalde stappen springt afhankelijk van de antwoorden die worden gekozen. \newline

%Omdat Power Automate een betalende licentie heeft, werd er eerst een tussenstap gemaakt. De gegevens werden naar een Excel gestuurd. In deze Excel werden de kolommen exact hetzelfde aangemaakt als in de test Excel voor het script. Deze zorgden voor een duidelijke verwerking van de gegevens. Op deze manier kon er ook eenvoudig gecontroleerd worden of de gebruikers de correcte gegevens meekregen. Dit was een tijdelijke, noodzakelijke tussenstap. \newline



%De volgende stap is het script op de server aanroepen. Dit kan je doen door een stap toe te voegen in de cloud. De stap heeft de naam ``Run a flow built with Power Automate for desktop''. Hierin kan je een flow kiezen, die je op de desktop gemaakt hebt. 

%Om dit te kunnen doen moet je uiteraard een desktop hebben om deze flow op te runnen. Hiervoor moet ``Power automate desktop'' geïnstalleerd worden. Dit kan gewoon via de website van Microsoft. Eens geïnstalleerd, moet je inloggen met het account waarmee je ook in de cloud bent ingelogd. Hierna kan je een flow maken. Je kan de flow uiteindelijk ook zien staan in de cloud, maar zal deze nooit kunnen bewerken in de cloud. Dit moet on-premesis gebeuren.

%Om een connectie te maken tussen de cloud versie van Power Automate en de desktop versie, moet je een andere applicatie gebruiken.  Hiervoor wordt de applicatie ``Power automate machine runtime'' gebruikt. Deze geeft meteen de correcte omgeving weer, waarin je aan het werken bent. Ook al geeft deze de correcte omgeving weer, wilt dit niet zeggen dat er al een connectie is. Hiervoor moet je op de knop ``Refresh'' drukken. Hierdoor maakt de machine dan connectie met het internet en met jouw cloud. 

%Eens dit is gedaan, kan je aan de flow beginnen. De flow in de desktop versie is redelijk simpel. Deze roept het script, dat lokaal opgeslagen staat op dezelfde computer of server, aan. Hiervoor wordt de stap ``Run PowerShell Script'' gebruikt. In deze stap wordt het script aangeroepen vanop de locatie waar het staat. Bij het script horen ook variabelen. Deze dienen ook aangemaakt te worden in de desktop applicatie. Deze variabelen hebben exact dezelfde naam als de gegevens die worden toegewezen in de cloud versie. Op deze manier kunnen ze aan elkaar gelinkt worden.

%De volgende stappen in de desktop versie, zijn enkel om de correcte gebruikersnaam door te geven aan de cloud. Dit is nodig omdat de gebruiker in de cloud ook toegevoegd wordt aan groepen. Hiervoor is de correcte naam nodig, deze wordt dan ook vanuit de desktop versie doorgegeven aan de cloud versie om de foutmarge zo klein mogelijk te maken.

%Verder is er nog 1 stap die in de desktop versie moet gebeuren. Deze stap zorgt ervoor dat de gebruiker bij tijdelijke contracten ook een einddatum krijgen in de Active Directory. Om veiligheidsredenen wordt dit nu handmatig gedaan als er een gekende einddatum is, maar bij een automatisch proces, gebeurt dat het beste automatisch. Deze stap zit niet in het script verwerkt aangezien de code om dit te doen, niet werkt in het script maar wel daar buitenom. Een verklaring hiervoor is niet gevonden. \newline

%Hierna moet dan nog de laatste stappen gebeuren voor de nieuwe gebruiker. Het toekennen van licentiesin entra id. Dit gebeurt door de licenties toe te kennen aan groepen en gebruikers dan in deze groepen te plaatsen. Voordat dit kan gebeuren, moet de gebruiker aangemaakt worden in de cloud. Dit kan eenvoudig gebeuren door eerst een variabele aan te maken. Dit hoort een boolean te zijn. Deze waarde moet op false staan. 

%Hierna komt de grootste stap. Om de loop te maken moet de stap ``Do Until'' aangeroepen worden. Deze stap werkt eigenlijk zoals een do while in andere programmeertalen. Deze stap gaat controleren of de variabele van waarde is veranderd. Van false naar true dus. Binnen deze stap, zitten verschillende andere stappen. Zo wordt er eerst gekeken of de gebruiker aangemaakt is in ``Entra ID'', de opvolger van ``Azure AD''. De stap ``Get user profile (V2)'' wordt hiervoor gebruikt. Deze haalt gegevens op uit Entra ID. In de flow wordt dan de mail van de nieuwe gebruiker meegegeven. Als deze bestaat, zullen de gegevens in de output van dit object komen.Hierna wordt er een condition aangeroepen als stap. Deze stap gaat kijken of er een object van het user profiel uit de vorige stap, bestaat. Als deze niet leeg is, dan wilt het dus zeggen dat hij aangemaakt is. Hierna wordt de variabele aangepast van false naar true en zal de ``Do Until'' eindigen en zullen de volgende stappen uitgevoerd worden. 

%Stel dat de user nog niet zou bestaan, gaat de ``Get user profile'' stap een foutmelding geven. Deze error kan je omzeilen door te zeggen dat de condition moet blijven lopen, ook als de stap ervoor een foutmelding geeft. 

%Verder moet er nog een ``Delay'' stap komen. Deze stap zorgt dat de flow niet iedere seconde een vraag stuurt of deze al bestaat. Mocht dit wel zo zijn, zou er een te grote hoeveelheid aan vragen zijn en zou de flow blokkeren. De vertraging staat ingesteld op iedere 30 seconden en zal 120 keer lopen. Dit om zeker te zijn dat de automatische synchronisatie cyclus zeker voldoende tijd heeft om te lopen. Deze loopt 1 keer per half uur. \newline

%De laatste stappen zijn dan om de gebruiker toe te voegen aan de correcte groepen in Entra. Dit zijn groepen die licenties meegeven. Dit gebeurt met de stap ``Office 365 Groups''. Deze worden toegevoegd op vlak van het group-id in Entra en de User Principal Name. \newline

%Eventueel kan er daarna nog een stap komen om terug te koppelen naar het algemene systeem voor de onboarding. Momenteel kan dit nog niet toegepast worden, aangezien hiervoor nog geen systeem bestaat en dit ook nog in productie zit. \newline

%Figuur \ref{fig:PowerAutomateCloud} toont een visuele voorstelling van de PowerAutomate flow.

%\begin{figure}[htbp]
%\centering
%\includegraphics[width=0.8\textwidth]{PowerAutomate.png}
%\caption[Complete Power Automate Cloud flow]{Deze afbeelding is de complete flow die in de cloud is. Deze flow heeft alle tussenstappen die hierboven beschreven zijn en waar er duidelijk vermeld staat dat het een cloud stap is. }
%\label{fig:PowerAutomateCloud}
%\end{figure}
